\section{Evidence and formalization boundary}
\label{app:boundary}

We distinguish four kinds of evidence and label every quantitative claim by the
kind that supports it. These are different kinds of evidence, not
interchangeable levels of confidence; in particular, compiled algebra does not
formalize the surrounding stochastic theorem.

\begin{itemize}
\item[\textbf{C}] Conventional mathematical proof. This covers the controlled
construction of Section~\ref{sec:model}, the segment lemma, the exposure floor
(Theorem~\ref{thm:main}), the actuator corollary, the amplitude floor
(Theorem~\ref{thm:amp}), the time floor (Proposition~\ref{prop:time}), the
weighted-drift proposition, the logarithmic order (Corollary~\ref{cor:order}),
the dimension-free sufficiency of Proposition~\ref{prop:kappa}, and the
concentration and administered-amount bounds of Section~\ref{sec:pkpd}.
\item[\textbf{E}] Exact finite computation or rigorous enclosure. This covers
all rational source residuals of Sections~\ref{sec:source} and
\ref{sec:memory}, every entry of Tables~\ref{tab:sites}, \ref{tab:amp} and
\ref{tab:expN}, the determinant polynomial \eqref{eq:detpoly}, the Taylor
coefficient of Appendix~\ref{app:taylor}, and the outward-rounded dyadic
enclosure \eqref{eq:validated}.
\item[\textbf{K}] Compiled declarations in Lean 4 with mathlib
\citep{moura2021,mathlib2020}, of the stated scope only: the literal six-state
source algebra, the supersolution segment identity, the control inequality
\eqref{eq:slacks}, the improved weight, the $RR$ coefficient calculation, and
auxiliary finite barrier and dose rearrangement identities. Strict receipts
record the pinned toolchain, source hashes and warnings-as-errors verification.
\item[\textbf{N}] Exploratory floating computation. This covers the schedule
curves and the full-withdrawal comparison of Figure~\ref{fig:policy}, the
large-$N$ values of Remark~\ref{rem:bigN} and the open markers of
Figure~\ref{fig:sites}, the candidate vectors later verified exactly, and the
floating amplitude optimization quoted in Section~\ref{sec:expN}.
\end{itemize}

Neither the main stochastic theorems nor the validated integrator is claimed to
be formalized in a proof assistant. That boundary is part of the result and not
an omitted assumption hidden in an axiom. The amplitude floors of
Table~\ref{tab:amp} are a natural candidate for a further finite Lean module,
since they reduce to two systems of rational polynomial inequalities in a fixed
vector; the stochastic step that turns them into
Proposition~\ref{prop:ampfloor} is Theorem~\ref{thm:amp} and is conventional.

The accompanying source package contains \texttt{main.tex} and the section
files, \texttt{references.bib}, vector figures, the certificate data files, the
generators and checkers named in Appendix~\ref{app:certs}, and a README with
canonical commands. The companion source manuscript for the inheritance
construction is \citet{molecular}.
